\chapter{Datos y universo de inversión}
\label{cap:datos}

Este capítulo describe de dónde salen los datos, cómo se define el universo de acciones sobre el
que opera el sistema y qué sesgos introduce cada decisión de origen. La tesis del capítulo es que
gran parte del rigor de un estudio sobre mercados se juega \emph{antes} de entrenar ningún modelo,
en cómo se construye la base de datos: un universo mal definido o una fuente con sesgos ocultos
contamina cualquier resultado posterior por sofisticado que sea el algoritmo.

\section{Fuentes de datos}

El sistema combina tres fuentes gratuitas, cada una elegida por lo que aporta y con conciencia de
sus límites:

\begin{itemize}
  \item \textbf{Finnhub} proporciona los datos fundamentales de las empresas (magnitudes contables
  y ratios). Es la fuente de las variables con las que se construyen los factores de calidad,
  crecimiento y valoración.

  \item \textbf{Yahoo Finance} proporciona los precios históricos ajustados de las acciones y del
  índice de referencia, base de los factores de \emph{momentum} y del cálculo de retornos y del
  \emph{backtest}.

  \item \textbf{SEC EDGAR} proporciona un dato que las otras dos no dan y que es crítico para la
  metodología: la \textbf{fecha real de presentación} de cada informe (\texttt{filingDate}), es
  decir, el día en que un dato fundamental pasó a ser público. Es gratuita y cubre desde 1993. Su
  papel en la eliminación del \emph{lookahead bias} se desarrolla en el
  capítulo~\ref{cap:metodologia}.
\end{itemize}

La combinación no es redundante: Finnhub dice \emph{cuánto} valía una magnitud, EDGAR dice
\emph{cuándo} se supo, y Yahoo aporta el precio con el que esa información se convierte en retorno.

\section{Universo dinámico del S\&P 500}

El universo de acciones no es la composición \emph{actual} del S\&P 500, sino su \textbf{composición
histórica real por fecha}. Una señal generada en 2016 solo puede ver las empresas que formaban
parte del índice en 2016; la lógica de pertenencia por fecha está encapsulada en
\texttt{module/data/universe.py} y se consulta en cada fecha de decisión.

Esta decisión elimina el \textbf{sesgo de inclusión anticipada}: si se usara la composición actual,
el sistema podría comprar en 2010 una empresa que solo entró en el índice en 2018 precisamente
porque le fue bien, incorporando información del futuro. Usar la composición vigente en cada fecha
cierra esa vía. Conviene ser honesto sobre lo que este control \emph{no} resuelve: restringir el
universo al S\&P 500 sigue siendo una decisión ---son grandes compañías, no el mercado entero---,
pero es una restricción declarada y estable, no un sesgo silencioso.

\section{Reciclaje de \emph{tickers}}

Un problema sutil de trabajar con series largas es que un mismo símbolo de cotización puede haber
identificado a empresas distintas en el tiempo: un \emph{ticker} se libera cuando una empresa
desaparece y puede reasignarse años después a otra. Sin una guarda, el sistema encadenaría los
precios de dos compañías distintas como si fueran una, generando retornos ficticios en el empalme.
El proyecto incorpora una guarda de reciclaje de \emph{tickers} (\texttt{module/data/universe.py}) que
impide mezclar identidades bajo un mismo símbolo; casos reales como los símbolos \texttt{CPQ} o
\texttt{MOB} ilustran por qué la comprobación es necesaria y no meramente teórica.

\section{Sesgo de supervivencia, medido por año}

El sesgo de supervivencia es probablemente la limitación más importante del origen de datos, y el
proyecto opta por \textbf{medirlo} en lugar de limitarse a declararlo. Las fuentes gratuitas no
conservan datos fundamentales de las empresas que quebraron o desaparecieron, de modo que la
cobertura de datos es desigual a lo largo del tiempo: parte de valores del orden del 50\,\% de las
compañías del índice a comienzos de la década de 2000 y sube hasta cifras próximas al 90\,\% en los
años más recientes. Esta cobertura por año es un artefacto verificable del proyecto, no una
estimación cualitativa.

La consecuencia metodológica es directa: el sistema se \textbf{centra en el periodo 2016 en
adelante}, donde la cobertura es más alta y el sesgo de supervivencia, aunque no nulo, es
sensiblemente menor. Es una decisión declarada y basada en el diagnóstico de cobertura, no una
elección oportunista de la ventana que mejor resultado da; su impacto y sus límites se retoman en
el capítulo~\ref{cap:limitaciones}.
